%==============================================================================
%  MATH 231 -- Linear Algebra I (Fall 2026)
%  A Natural Progression Through the Course Material
%
%  This document re-sequences the tentative syllabus for dependency-correctness
%  WITHOUT omitting any proposed material.  Each unit specifies:
%     - Domain(s)          : Pure LA / Numerical LA / Optimization / ML /
%                            Multivariable Calculus / Spectral Graph Theory / ...
%     - Primary object(s)  : vectors, matrices, complex numbers, quaternions, sets
%     - Topics             : the conceptual content
%     - Key theorems       : named results to prove or state
%     - Tools / formulas   : closed-form computational instruments
%     - Algorithms         : procedures to implement
%     - Applications        : where it is used (ML, physics, DSP, ...)
%     - ENTER              : prerequisites needed to begin the unit
%     - EXIT (gateway)     : what a student must master to advance
%
%  Priority tags:  [Core]  [Recommended]  [Enrichment]
%  Compile with pdflatex (standard TeX Live packages only).
%==============================================================================
\documentclass[11pt]{article}

\usepackage[margin=1in]{geometry}
\usepackage{amsmath,amssymb,amsthm}
\usepackage[dvipsnames]{xcolor}
\usepackage{enumitem}
\usepackage{booktabs}
\usepackage{array}
\usepackage[hidelinks]{hyperref}
\usepackage{titlesec}

%------------------------------------------------------------------ style ----
\titleformat{\section}{\Large\bfseries\sffamily}{\thesection.}{0.6em}{}
\titleformat{\part}{\centering\huge\bfseries\sffamily}{}{0pt}{}

\newcommand{\dom}[1]{{\sffamily\bfseries\color{RoyalBlue}#1}}
\newcommand{\pri}[1]{{\footnotesize\sffamily\color{BrickRed}[\,#1\,]}}
\newcommand{\obj}[1]{{\itshape #1}}

% Per-unit field list -------------------------------------------------------
\newenvironment{unitfields}
  {\begin{description}[leftmargin=0pt,labelindent=0pt,style=nextline,
      font=\normalfont\sffamily\bfseries\color{RoyalBlue},itemsep=3pt,topsep=4pt]}
  {\end{description}}

\newcommand{\ff}[1]{\item[#1]}

\theoremstyle{definition}
\newtheorem*{gateway}{Gateway to next unit}

\title{\vspace{-1.2cm}\textbf{\sffamily MATH 231 -- Linear Algebra I}\\[4pt]
       {\large\sffamily A Natural Progression Through the Course Material}\\[2pt]
       {\normalsize\sffamily Fall 2026 -- working draft}}
\author{}
\date{\normalsize\sffamily Draft dated July 1, 2026}

\begin{document}
\maketitle
\vspace{-1.2cm}

%==============================================================================
\section*{How to read this document}
%==============================================================================
The tentative syllabus contains \emph{more than one semester} of material by
design: the intent is a wide breadth from which to select the topics students
find most compelling and the instructor deems most essential.  Accordingly:

\begin{itemize}[leftmargin=1.4em,itemsep=2pt]
  \item Units are ordered so that every concept is preceded by its prerequisites.
        The \dom{Domain} tag names the discipline; the \obj{objects} name the
        primary mathematical objects in play.
  \item Priority tags let you carve a one-semester path:
        \pri{Core} = spine of the course (skip nothing);
        \pri{Recommended} = strongly desired if time permits;
        \pri{Enrichment} = optional depth / seminar-style topics.
  \item Each unit ends with an explicit \textbf{Gateway}: the competencies a
        student must hold before the next unit is intelligible.
  \item The course is grouped into four \emph{arcs}. A minimal one-semester
        realization is Arcs~I--III plus a chosen slice of Arc~IV.
\end{itemize}

\noindent\textbf{Domains used:} Pure Linear Algebra; Numerical Analysis /
Numerical Linear Algebra; Multivariable Calculus; Optimization; Machine
Learning; Spectral Graph Theory.\\[2pt]
\textbf{Implementation language:} Python.\\[4pt]
\textbf{Textbook reference keys} (each unit lists chapter/lecture-level pointers,
not page-exact):
\textbf{[BV]} Boyd \& Vandenberghe, \emph{Introduction to Applied Linear Algebra};
\textbf{[TB]} Trefethen \& Bau, \emph{Numerical Linear Algebra} (cited by
\emph{Lecture});
\textbf{[S]} Strang, \emph{Linear Algebra and Its Applications} (4th ed.);
\textbf{[A]} Axler, \emph{Linear Algebra Done Right} (4th ed.).
Topics lying outside all four texts are marked \emph{supplementary}.

%==============================================================================
\section*{Dependency map (arc-level)}
%==============================================================================
\begin{center}
\renewcommand{\arraystretch}{1.25}
\begin{tabular}{@{}c l l@{}}
\toprule
\textbf{Arc} & \textbf{Units} & \textbf{Depends on} \\
\midrule
I.   Foundations                 & 1--3   & --- \\
II.  Matrices \& Continuous Math  & 4--6   & Arc I \\
III. Core Matrix Theory          & 7--9   & Arcs I--II \\
IV.  Advanced Methods \& Apps     & 10--12 & Arcs I--III \\
\bottomrule
\end{tabular}
\end{center}

\noindent Two ``honest'' forward references worth naming up front:
the \emph{spectral} and \emph{nuclear} matrix norms are \emph{defined} in
Unit~4 but only \emph{computed/justified} after the SVD in Unit~11; the
\emph{Hessian} is introduced in Unit~5 but its definiteness analysis is
completed in Unit~11.

%##############################################################################
\part*{Arc I -- Foundations}
%##############################################################################

%------------------------------------------------------------------ Unit 1 ---
\section{Vectors, Vector Spaces \& the Language of the Course}
\begin{unitfields}
\ff{Domain(s)} \dom{Pure Linear Algebra} \pri{Core}
\ff{Primary object(s)} \obj{sets, fields ($\mathbb{R},\mathbb{C}$), vectors}
\ff{Topics}
  Sets and functions; the field axioms for $\mathbb{R}$ and $\mathbb{C}$.
  Vectors in $\mathbb{R}^2,\mathbb{R}^3,\mathbb{R}^n$ and higher-dimensional
  spaces. The algebraic properties: vector addition, scalar multiplication,
  additive identity/inverse, commutativity, associativity, the zero vector,
  distributivity, multiplicative identity. \textbf{Formal definition of a
  vector space} (a set $V$ with $+$ and scalar multiplication satisfying the
  six axioms). Subspaces; span; linear independence/dependence; basis;
  dimension. Geometric picture: parallelogram rule, scaling, normalization.
\ff{Key theorems}
  Uniqueness of the additive identity and additive inverse;
  every spanning set contains a basis / every independent set extends to a
  basis \pri{Recommended}; dimension is well-defined (basis-size invariance)
  \pri{Recommended}.
\ff{Tools / formulas}
  Linear combinations; coordinate representation relative to a basis.
\ff{Algorithms} --- (algorithmic content begins in Unit~2).
\ff{Applications}
  Image vectorization (an image as a point in $\mathbb{R}^{HW}$);
  physical vectors (position/velocity/acceleration) as a motivating instance
  \pri{Recommended}.
\ff{Languages} Python (construct/visualize vectors).
\ff{Textbook refs}
  \textbf{[A]} Ch.~1--2 (vector spaces, independence, bases, dimension);
  \textbf{[S]} Ch.~2; \textbf{[BV]} Ch.~1, 5.
\ff{ENTER} High-school algebra; comfort with set notation.
\end{unitfields}
\begin{gateway}
Student can state the six vector-space axioms, verify whether a given set with
operations is a vector space, and compute a basis/dimension of a subspace of
$\mathbb{R}^n$.
\end{gateway}

%------------------------------------------------------------------ Unit 2 ---
\section{Inner Products, Norms \& Similarity}
\begin{unitfields}
\ff{Domain(s)} \dom{Pure Linear Algebra} + \dom{Numerical LA} + \dom{Machine Learning} \pri{Core}
\ff{Primary object(s)} \obj{vectors}
\ff{Topics}
  Dot (inner) product; component and orthogonal projection of one vector onto
  another; angle between vectors. Orthogonal vs.\ parallel vectors.
  \textbf{Vector norms}: Euclidean/standard ($\ell_2$), Manhattan/taxicab
  ($\ell_1$), Chebyshev/chessboard ($\ell_\infty$), and the general $p$-norm.
  Normalization. Distance/metric induced by a norm. Cosine similarity and
  general similarity kernels.
\ff{Key theorems}
  \textbf{Cauchy--Schwarz inequality}; triangle inequality; Pythagorean
  identity for orthogonal vectors \pri{Recommended}.
\ff{Tools / formulas}
  Dot product $\langle x,y\rangle$; projection
  $\operatorname{proj}_y x=\frac{\langle x,y\rangle}{\langle y,y\rangle}y$;
  cosine similarity $\frac{\langle x,y\rangle}{\|x\|\,\|y\|}$; $p$-norm
  $\|x\|_p=(\sum|x_i|^p)^{1/p}$; general kernel $k(x,y)$.
\ff{Algorithms}
  \textbf{Gram--Schmidt orthogonalization} (classical form; modified form
  deferred to Unit~10) \pri{Core};
  \textbf{$k$-Nearest Neighbors} \pri{Recommended};
  \textbf{$k$-means clustering} \pri{Recommended}.
\ff{Applications}
  Semantic embeddings in language models (cosine similarity between token/word
  vectors); retrieval; clustering; a geometric preview of the SVM margin
  \pri{Enrichment}.
\ff{Languages}
  Python (inner product, cosine similarity, KNN, $k$-means) \pri{Core}.
\ff{Textbook refs}
  \textbf{[A]} Ch.~6 (inner product spaces);
  \textbf{[TB]} Lec.~2--3 (orthogonal vectors, norms), Lec.~8 (Gram--Schmidt);
  \textbf{[S]} Ch.~3; \textbf{[BV]} Ch.~3--5.
\ff{ENTER} Unit~1 (vectors, span, basis).
\end{unitfields}
\begin{gateway}
Student can compute norms/projections/cosine similarity by hand and in code,
prove Cauchy--Schwarz, and run Gram--Schmidt to produce an orthonormal set.
\end{gateway}

%------------------------------------------------------------------ Unit 3 ---
\section{Complex Numbers \& Quaternions}
\begin{unitfields}
\ff{Domain(s)} \dom{Pure Linear Algebra / Algebra} \pri{Core (complex)}, \pri{Enrichment (quaternions)}
\ff{Primary object(s)} \obj{complex numbers $\mathbb{C}$, quaternions $\mathbb{H}$}
\ff{Topics}
  Algebraic structure of $\mathbb{C}$: addition, multiplication, conjugation.
  Real and imaginary parts; the modulus. Polar form; the complex plane and the
  geometric interpretation of multiplication as rotation-and-scaling.
  Quaternions: non-commutative multiplication, use in 3-D rotation
  \pri{Enrichment}.
\ff{Key theorems}
  \textbf{Euler's formula} $e^{i\theta}=\cos\theta+i\sin\theta$;
  \textbf{De Moivre's theorem}; Fundamental Theorem of Algebra (stated)
  \pri{Recommended}.
\ff{Tools / formulas}
  Modulus $|z|=\sqrt{a^2+b^2}$; $\operatorname{Re}z,\operatorname{Im}z$;
  polar form $z=re^{i\theta}$; $n$-th roots of unity \pri{Recommended};
  quaternion product / rotation formula \pri{Enrichment}.
\ff{Algorithms} Complex arithmetic; rotation composition via quaternions \pri{Enrichment}.
\ff{Applications}
  Signal representation (phasors) \pri{Recommended}; 3-D rotations in graphics
  and robotics \pri{Enrichment}; foreshadowing complex eigenvalues (Unit~11).
\ff{Languages} Python (complex plane visualization).
\ff{Textbook refs}
  \textbf{[A]} Ch.~1 ($\mathbb{F}=\mathbb{R}$ or $\mathbb{C}$),
  Ch.~4 (polynomials, fundamental theorem of algebra);
  Euler's formula, De Moivre, quaternions \emph{supplementary}.
\ff{ENTER} Unit~1 (fields, vectors); Unit~2 (norms $\Rightarrow$ modulus).
\end{unitfields}
\begin{gateway}
Student can move fluently between rectangular and polar form, apply Euler/De
Moivre, and interpret complex multiplication geometrically.
\end{gateway}

%##############################################################################
\part*{Arc II -- Matrices \& Continuous Math}
%##############################################################################

%------------------------------------------------------------------ Unit 4 ---
\section{Matrices: Algebra, Determinant, Trace \& Matrix Norms}
\begin{unitfields}
\ff{Domain(s)} \dom{Pure Linear Algebra} + \dom{Numerical LA} \pri{Core}
\ff{Primary object(s)} \obj{matrices}
\ff{Topics}
  Matrices as arrays and as objects: addition, scalar multiplication, matrix
  multiplication, transpose/conjugate-transpose, inverse. Special matrices
  (identity, diagonal, symmetric, orthogonal). Determinant and its properties;
  trace. The matrix exponential $e^{A}$ (defined by its series; closed forms
  deferred to Unit~11). \textbf{Matrix norms}: induced/operator $1$-norm,
  $\infty$-norm, spectral ($2$-)norm; Frobenius norm; nuclear norm.
\ff{Key theorems}
  Determinant multiplicativity $\det(AB)=\det A\det B$; invertibility
  $\iff\det\neq 0$; submultiplicativity of induced norms
  $\|AB\|\le\|A\|\|B\|$ \pri{Recommended}.
\ff{Tools / formulas}
  $\operatorname{tr}(A)=\sum a_{ii}$; determinant (cofactor and via
  elimination); $\|A\|_F=\sqrt{\sum a_{ij}^2}$; induced norm
  $\|A\|=\sup_{x\neq0}\|Ax\|/\|x\|$; $e^{A}=\sum_k A^k/k!$.
  \emph{Note:} spectral and nuclear norms are \emph{defined} here but
  \emph{computed} via the SVD in Unit~11.
\ff{Algorithms}
  Naive and blocked matrix multiply; matrix--vector product \pri{Core}.
\ff{Applications}
  Linear operators on data; adjacency/weight matrices (preview of graphs).
\ff{Languages}
  Python (matmul, mat--vec, blocking and timing experiments) \pri{Core}.
\ff{Textbook refs}
  \textbf{[S]} Ch.~1 (matrix operations), Ch.~4 (determinants);
  \textbf{[A]} Ch.~3, Ch.~9 (trace, determinant);
  \textbf{[TB]} Lec.~1, 3; \textbf{[BV]} Ch.~6, 10--11;
  matrix exponential \emph{supplementary}.
\ff{ENTER} Units~1--2 (vectors, norms).
\end{unitfields}
\begin{gateway}
Student can multiply/invert matrices, compute determinant and trace, and
evaluate the Frobenius and induced $1/\infty$ norms; understands what the
spectral/nuclear norms \emph{mean} even before computing them.
\end{gateway}

%------------------------------------------------------------------ Unit 5 ---
\section{Multivariable Calculus: Gradients, Hessians \& Taylor's Theorem}
\begin{unitfields}
\ff{Domain(s)} \dom{Multivariable Calculus} \pri{Core}
\ff{Primary object(s)} \obj{multivariate functions, vectors, matrices}
\ff{Topics}
  Multivariable scalar functions and vector-valued functions; their graphs.
  Equations of planes; tangent planes. The definition of the derivative for
  multivariate scalar functions; partial derivatives. Linear approximation.
  \textbf{The gradient} (a vector) and \textbf{the Hessian} (a matrix;
  definiteness analysis completed in Unit~11). Review of exponential and
  logarithmic functions and \emph{derivations} of the derivatives of $e^x$ and
  $\ln x$.
\ff{Key theorems}
  \textbf{Taylor's theorem}, univariate and multivariate (first- and
  second-order forms) \pri{Core}; equality of mixed partials (Clairaut)
  \pri{Recommended}.
\ff{Tools / formulas}
  Gradient $\nabla f$; Hessian $H_f$; tangent-plane / first-order model
  $f(x)\approx f(a)+\nabla f(a)^\top(x-a)$; second-order model with
  $\tfrac12(x-a)^\top H(x-a)$; $\frac{d}{dx}e^x=e^x$,
  $\frac{d}{dx}\ln x=1/x$.
\ff{Algorithms}
  Numerical gradient (finite differences) \pri{Recommended}.
\ff{Applications}
  Sets the analytic foundation for all of Arc~IV (optimization, regression).
\ff{Languages}
  Python (surface/contour plots, gradient fields).
\ff{Textbook refs}
  \textbf{[BV]} App.~C (derivatives \& optimization);
  Taylor's theorem and the Hessian \emph{supplementary}
  (multivariable-calculus text).
\ff{ENTER} Units~1--2 (vectors, inner products) and Unit~4 (matrices, for the Hessian).
\end{unitfields}
\begin{gateway}
Student can compute gradients/Hessians, write first- and second-order Taylor
models of a multivariate function, and derive $\frac{d}{dx}e^x$ and
$\frac{d}{dx}\ln x$ from first principles.
\end{gateway}

%------------------------------------------------------------------ Unit 6 ---
\section{Numerical Foundations: Complexity, Floating Point, Conditioning \& Stability}
\begin{unitfields}
\ff{Domain(s)} \dom{Numerical Analysis} \pri{Core (short unit)}
\ff{Primary object(s)} \obj{real numbers, floating-point numbers, functions}
\ff{Topics}
  Time complexity: Big-$O$, $\Omega$, $\Theta$. The floating-point number
  system as the footing of numerical analysis. The \textbf{condition number} of
  a (univariate, then matrix) problem $\Rightarrow$ sensitivity of outputs to
  input perturbation. \textbf{Stability} of algorithms. Sources of numerical
  error: cancellation, rounding, precision. Tolerances and stopping criteria
  in iterative algorithms.
\ff{Key theorems}
  Relationship $\text{error}\lesssim(\text{condition number})\times
  (\text{machine epsilon})$ for stable algorithms \pri{Recommended}.
\ff{Tools / formulas}
  Machine epsilon; condition number $\kappa(f)=|xf'(x)/f(x)|$ (univariate) and
  $\kappa(A)=\|A\|\,\|A^{-1}\|$ (matrix, referenced forward to Unit~11);
  absolute vs.\ relative error.
\ff{Algorithms}
  Illustrative unstable vs.\ stable computations (e.g., catastrophic
  cancellation demos) \pri{Core}.
\ff{Applications}
  Governs the reliability of every later algorithm (elimination, QR, eigen).
\ff{Languages}
  Python (float vs.\ double behavior, machine epsilon,
  ill-conditioned examples).
\ff{Textbook refs}
  \textbf{[TB]} Lec.~12--15 (conditioning, floating point, stability);
  Big-Oh / complexity \emph{supplementary}.
\ff{ENTER} Unit~4 (matrix norms); Unit~5 (derivatives, for condition numbers).
\end{unitfields}
\begin{gateway}
Student can give asymptotic costs, explain what conditioning and stability mean,
identify cancellation, and set a sensible tolerance for an iterative method.
\end{gateway}

%##############################################################################
\part*{Arc III -- Core Matrix Theory}
%##############################################################################

%------------------------------------------------------------------ Unit 7 ---
\section{Linear Systems \& Elimination}
\begin{unitfields}
\ff{Domain(s)} \dom{Pure Linear Algebra} + \dom{Numerical LA} \pri{Core}
\ff{Primary object(s)} \obj{matrices, vectors}
\ff{Topics}
  Systems $Ax=b$; row operations; row-echelon and reduced row-echelon form.
  Existence and uniqueness of solutions; rank; pivots; homogeneous vs.\
  inhomogeneous systems. \textbf{LU factorization}; pivoting for stability
  (tie to Unit~6).
\ff{Key theorems}
  Rouch\'e--Capelli (solvability via rank) \pri{Recommended};
  existence/uniqueness of $LU$ with partial pivoting \pri{Recommended}.
\ff{Tools / formulas}
  Gaussian/Gauss--Jordan elimination; back/forward substitution; $A=LU$,
  $PA=LU$.
\ff{Algorithms}
  \textbf{Gaussian elimination}; \textbf{LU factorization} with partial
  pivoting \pri{Core}.
\ff{Applications}
  The computational workhorse behind regression, Newton steps, and more.
\ff{Languages}
  Python (elimination by hand, \texttt{scipy.linalg.lu} verification)
  \pri{Core}.
\ff{Textbook refs}
  \textbf{[S]} Ch.~1 (Gaussian elimination);
  \textbf{[TB]} Lec.~20--22 (elimination, pivoting, stability);
  \textbf{[BV]} Ch.~8, 11.
\ff{ENTER} Units~1--2, 4 (vectors, matrices); Unit~6 (why pivoting).
\end{unitfields}
\begin{gateway}
Student can solve $Ax=b$ by elimination, determine solvability from rank, and
produce/use an $LU$ factorization.
\end{gateway}

%------------------------------------------------------------------ Unit 8 ---
\section{Linear Transformations}
\begin{unitfields}
\ff{Domain(s)} \dom{Pure Linear Algebra} \pri{Core}
\ff{Primary object(s)} \obj{linear maps, matrices}
\ff{Topics}
  Linear maps $T:V\to W$; the matrix of a linear map; composition as matrix
  product. Kernel (null space) and image (range) of a map. Injectivity/
  surjectivity/invertibility. Change of basis and similarity.
\ff{Key theorems}
  \textbf{Rank--Nullity theorem}; a linear map is invertible $\iff$ its matrix
  is invertible \pri{Core}; change-of-basis / similarity relation
  $B=P^{-1}AP$ \pri{Recommended}.
\ff{Tools / formulas}
  Matrix representation in a basis; change-of-basis matrix $P$.
\ff{Algorithms}
  Computing kernel/image bases via elimination \pri{Core}.
\ff{Applications}
  Rotations/reflections/projections as maps; graphics transforms
  \pri{Recommended}; sets up the four subspaces and eigen-theory.
\ff{Languages} Python (visualize $2$-D/$3$-D transforms).
\ff{Textbook refs}
  \textbf{[A]} Ch.~3 (linear maps);
  \textbf{[S]} App.~A (linear transformations);
  \textbf{[BV]} Ch.~2 (linear functions).
\ff{ENTER} Unit~7 (elimination, rank); Units~1--2.
\end{unitfields}
\begin{gateway}
Student can pass between a linear map and its matrix, compute kernel/image,
apply rank--nullity, and change basis.
\end{gateway}

%------------------------------------------------------------------ Unit 9 ---
\section{The Four Fundamental Subspaces}
\begin{unitfields}
\ff{Domain(s)} \dom{Pure Linear Algebra} \pri{Core}
\ff{Primary object(s)} \obj{subspaces, matrices}
\ff{Topics}
  Column space, null space, row space, left null space of $A$; their
  dimensions and how they follow from rank. Orthogonality relations between the
  pairs. The geometry that underlies projection and least squares.
\ff{Key theorems}
  \textbf{Fundamental Theorem of Linear Algebra} (dimensions + orthogonality:
  $\mathcal{N}(A)\perp\mathcal{R}(A^\top)$, $\mathcal{N}(A^\top)\perp
  \mathcal{R}(A)$) \pri{Core}.
\ff{Tools / formulas}
  $\dim\mathcal{C}(A)=\operatorname{rank}A=r$; $\dim\mathcal{N}(A)=n-r$;
  bases of each subspace from RREF.
\ff{Algorithms}
  Extract bases of all four subspaces from elimination \pri{Core}.
\ff{Applications}
  Direct foundation for least squares (Unit~10) and PCA/SVD geometry (Unit~11).
\ff{Languages} Python (rank, null-space bases).
\ff{Textbook refs}
  \textbf{[S]} Ch.~2 (\S2.4, the four subspaces);
  \textbf{[A]} Ch.~3 (null space \& range, fundamental theorem of linear maps);
  \textbf{[BV]} Ch.~5.
\ff{ENTER} Units~7--8 (elimination, maps, rank--nullity).
\end{unitfields}
\begin{gateway}
Student can name and produce bases for all four subspaces of a given $A$ and
state their dimensions and orthogonality relations.
\end{gateway}

%##############################################################################
\part*{Arc IV -- Advanced Methods \& Applications}
%##############################################################################

%------------------------------------------------------------------ Unit 10 --
\section{Orthogonality, Projections, Least Squares \& Factorizations}
\begin{unitfields}
\ff{Domain(s)} \dom{Numerical LA} + \dom{Machine Learning} \pri{Core}
\ff{Primary object(s)} \obj{matrices, vectors, subspaces}
\ff{Topics}
  Orthogonal projection onto a subspace; orthonormal bases. The least-squares
  problem $\min_x\|Ax-b\|_2$ and the normal equations $A^\top A x=A^\top b$.
  Linear regression as least squares. \textbf{QR factorization} (classical and
  \emph{modified} Gram--Schmidt; Householder reflections). \textbf{Cholesky
  factorization} of the SPD matrix $A^\top A$. Introduction to quadratic forms
  (definiteness deepened in Unit~11). Orthogonal transforms as a change to an
  orthonormal basis; the \textbf{Discrete Cosine Transform (DCT)} as a
  \emph{fixed} orthonormal basis, contrasted with the data-adaptive SVD basis
  (payoff in the JPEG/SVD image-compression demo, Unit~11) \pri{Recommended}.
\ff{Key theorems}
  Projection/orthogonality principle (residual $\perp$ column space);
  existence \& uniqueness of the least-squares solution when $A$ has full
  column rank \pri{Core}; $A=QR$ existence \pri{Recommended}.
\ff{Tools / formulas}
  Projection matrix $P=A(A^\top A)^{-1}A^\top$; normal equations;
  $A=QR$; $A^\top A=LL^\top$ (Cholesky); quadratic form $q(x)=x^\top Ax$.
\ff{Algorithms}
  \textbf{Ordinary Least Squares (OLS)} \pri{Core};
  \textbf{Modified Gram--Schmidt} \pri{Core};
  \textbf{Householder QR} \pri{Recommended};
  \textbf{Cholesky factorization} \pri{Recommended}.
\ff{Applications}
  Linear regression / curve fitting \pri{Core}; data whitening \pri{Enrichment}.
\ff{Languages}
  Python (linear regression demo; \texttt{numpy.linalg.qr} and
  \texttt{scipy.linalg.cholesky} verification) \pri{Core}.
\ff{Textbook refs}
  \textbf{[TB]} Lec.~6--11 (projectors, QR, Gram--Schmidt, Householder, least
  squares), Lec.~23 (Cholesky);
  \textbf{[BV]} Ch.~12--13; \textbf{[S]} Ch.~3, Ch.~6; \textbf{[A]} Ch.~6.
\ff{ENTER} Unit~9 (subspaces, projection geometry); Unit~6 (why MGS/Householder over normal equations).
\end{unitfields}
\begin{gateway}
Student can set up and solve a least-squares problem three ways (normal
equations, QR, Cholesky), fit a regression model, and explain why QR is more
stable than the normal equations.
\end{gateway}

%------------------------------------------------------------------ Unit 11 --
\section{Eigenvalues, Eigenvectors \& the Singular Value Decomposition}
\begin{unitfields}
\ff{Domain(s)} \dom{Pure LA} + \dom{Numerical LA} + \dom{Spectral Graph Theory} + \dom{Machine Learning} \pri{Core}
\ff{Primary object(s)} \obj{matrices, vectors, complex numbers}
\ff{Topics}
  Eigenvalues/eigenvectors; characteristic polynomial; connection to
  determinant ($\prod\lambda_i$) and trace ($\sum\lambda_i$). Diagonalization
  and similarity. Symmetric matrices and orthogonal diagonalization.
  \textbf{Quadratic forms and definiteness} (completes Unit~10). The
  \textbf{SVD} and its geometry; low-rank approximation; \emph{now} the
  spectral and nuclear norms of Unit~4 are computable. Why eigenvalue problems
  are solved \emph{iteratively}. Graph Laplacians and spectral clustering.
  Stochastic (Markov) matrices and the dominant-eigenvector interpretation that
  underlies PageRank.
\ff{Key theorems}
  \textbf{Spectral theorem} (symmetric/Hermitian $\Rightarrow$ orthonormal
  eigenbasis) \pri{Core};
  \textbf{Schur's theorem} (unitary triangularization) \pri{Recommended};
  \textbf{Singular Value Decomposition} \pri{Core};
  \textbf{Eckart--Young} (best low-rank approximation) \pri{Recommended};
  \textbf{Abel--Ruffini theorem} (no general radical solution for degree
  $\ge 5$ $\Rightarrow$ eigenvalues need iteration) \pri{Enrichment};
  \textbf{Perron--Frobenius theorem} (a nonnegative/stochastic matrix has a
  positive dominant eigenvalue with a nonnegative eigenvector; underlies
  PageRank and spectral clustering) \pri{Recommended}.
\ff{Tools / formulas}
  Characteristic polynomial $\det(A-\lambda I)=0$; $A=Q\Lambda Q^\top$;
  $A=U\Sigma V^\top$; $\|A\|_2=\sigma_{\max}$, $\|A\|_*=\sum\sigma_i$;
  $\kappa_2(A)=\sigma_{\max}/\sigma_{\min}$; matrix exponential via
  eigendecomposition $e^{A}=Qe^{\Lambda}Q^{-1}$; graph Laplacian $L=D-W$.
\ff{Algorithms}
  Power iteration; QR algorithm (stated) \pri{Recommended};
  \textbf{PCA} \pri{Core};
  \textbf{Randomized SVD} and \textbf{Hutchinson trace estimation}
  (\emph{Monte-Carlo application}) \pri{Enrichment};
  spectral clustering on graph Laplacians \pri{Recommended}.
\ff{Applications}
  PCA / dimensionality reduction; Eigenfaces / ``Eigencats'' \pri{Recommended};
  spectral graph theory and community detection \pri{Recommended};
  low-rank compression \pri{Enrichment}; Latent Semantic Analysis on a TF--IDF
  term--document matrix (document retrieval); PageRank; SVD/DCT image
  compression; static word-embedding geometry (see Projects appendix).
\ff{Languages}
  Python (PCA, SVD, Eigenfaces, spectral graph demos;
  \texttt{numpy.linalg.eig}/\texttt{svd} verification) \pri{Core}.
\ff{Textbook refs}
  \textbf{[A]} Ch.~5, Ch.~7 (spectral theorem, SVD), Ch.~8 (Schur);
  \textbf{[S]} Ch.~5--6; \textbf{[TB]} Lec.~4--5 (SVD),
  Lec.~24--31 (eigenvalue algorithms, computing the SVD);
  PageRank, Perron--Frobenius, spectral graph theory \emph{supplementary}.
\ff{ENTER} Units~8--10 (maps, subspaces, projections, quadratic forms); Unit~3 (complex eigenvalues); Unit~6 (why iterative).
\end{unitfields}
\begin{gateway}
Student can diagonalize symmetric matrices, compute/interpret an SVD, use it
for low-rank approximation and PCA, evaluate the spectral/nuclear norms and
$\kappa_2$, and build a graph Laplacian for clustering.
\end{gateway}

%------------------------------------------------------------------ Unit 12 --
\section{Numerical Optimization \& Machine-Learning Capstone}
\begin{unitfields}
\ff{Domain(s)} \dom{Optimization} + \dom{Machine Learning} \pri{Core (grad.\ descent, Newton); Recommended/Enrichment beyond}
\ff{Primary object(s)} \obj{multivariate functions, matrices, vectors}
\ff{Topics}
  Unconstrained minimization; descent directions; step sizes/line search;
  convergence and its dependence on conditioning (tie to Units~6,11).
  \textbf{Gradient descent}; \textbf{Newton's method} (Hessian solve $=$ a
  linear system, Unit~7). Nonlinear least squares: \textbf{Gauss--Newton} and
  \textbf{Levenberg--Marquardt}.
\ff{If time permits}
  The following are \emph{not} promised. If the semester allows, we will
  discuss: quasi-Newton \textbf{BFGS}; constrained optimization and the
  \textbf{support vector machine} under the Vapnik objective, with kernel
  methods reusing the Unit~2 kernels; stochastic gradient descent
  (\textbf{SGD}); and model-selection methodology (train/validation/test splits
  and cross-validation as estimates of generalization error). These are carried
  in \texttt{plan/} rather than committed to here.
\ff{Key theorems}
  First-/second-order optimality conditions (via Taylor, Unit~5);
  convexity $\Rightarrow$ global optimum \pri{Recommended};
  local quadratic convergence of Newton's method \pri{Enrichment}.
\ff{Tools / formulas}
  GD update $x_{k+1}=x_k-\alpha\nabla f$; Newton
  $x_{k+1}=x_k-H^{-1}\nabla f$; Gauss--Newton normal-equation step;
  LM damped step $(J^\top J+\lambda I)\Delta=-J^\top r$.
\ff{Algorithms}
  \textbf{Gradient descent} \pri{Core}; \textbf{Newton's method} \pri{Core};
  \textbf{Gauss--Newton} \pri{Recommended};
  \textbf{Levenberg--Marquardt} \pri{Recommended}.
\ff{Applications}
  Model fitting; the full pipeline tying the course together
  (embeddings $\to$ similarity $\to$ optimization); signal trilateration
  (GPS-style positioning).
\ff{Languages}
  Python (gradient descent, Newton, Gauss--Newton/LM demos) \pri{Core}.
\ff{Textbook refs}
  \textbf{[BV]} Ch.~18--19 (nonlinear \& constrained least squares:
  Gauss--Newton, Levenberg--Marquardt), App.~C;
  gradient descent and Newton \emph{supplementary}.
\ff{ENTER} Unit~5 (gradient/Hessian/Taylor); Unit~7 (linear solves); Unit~10 (least squares); Unit~11 (definiteness/conditioning).
\end{unitfields}
\begin{gateway}
Course capstone: student can derive and implement gradient descent and Newton's
method and choose Gauss--Newton/LM for nonlinear least squares, articulating how
each method's behavior depends on conditioning.
\end{gateway}

%==============================================================================
\section*{One-semester carving suggestion}
%==============================================================================
A realistic 14--15 week path keeps all \pri{Core} items and a chosen subset of
\pri{Recommended}:
\begin{itemize}[leftmargin=1.4em,itemsep=1pt]
  \item \textbf{Arc I} (Units 1--3): $\sim$3 weeks. Quaternions $\to$ Enrichment.
  \item \textbf{Arc II} (Units 4--6): $\sim$3 weeks. Matrix-exp $\to$ light.
  \item \textbf{Arc III} (Units 7--9): $\sim$3 weeks.
  \item \textbf{Arc IV} (Units 10--12): $\sim$4--5 weeks. Pick \emph{one} of
        \{Levenberg--Marquardt, spectral graph theory, randomized SVD\} as the
        enrichment capstone rather than all three.
\end{itemize}

\bigskip
\noindent\textbf{Reference texts.}
Boyd \& Vandenberghe, \emph{Introduction to Applied Linear Algebra} (applied
spine); Trefethen \& Bau, \emph{Numerical Linear Algebra} (Arc~II \& IV
numerics); Strang, \emph{Linear Algebra and Its Applications} (subspaces \&
intuition); Axler, \emph{Linear Algebra Done Right} (theoretical backbone,
Arcs~I \& III).

\end{document}
